% ==============================================================
% 摘要
% ==============================================================

\chapter*{摘\quad 要}
\addcontentsline{toc}{chapter}{摘要}

\textcolor{red}{[TODO: 全文完成后撰写中文摘要，约500字]}

面对中国日益严峻的人口老龄化与慢性病管理压力，社区2型糖尿病管理面临"人力有限、需求无限"的结构性矛盾。本文设计并实现了一个基于多智能体系统（Multi-Agent System, MAS）的社区老年糖尿病智能随访管理平台。系统以LangGraph为多智能体编排框架，以大语言模型（LLM）为智能增强引擎，采用"规则驱动+LLM能力增强"的混合架构，构建了包含患者数据Agent、风险分诊Agent、随访调度Agent、用药管理Agent和医生工作台Agent在内的五Agent协同体系。系统通过检索增强生成（RAG）技术整合《中国糖尿病防治指南（2024版）》等权威知识，实现了健康数据智能采集、三色风险预警、动态随访调度、用药依从管理等核心功能。前端采用适老化设计，支持语音交互，降低老年患者使用门槛。

\vspace{1em}
\noindent\textbf{关键词：}多智能体系统；大语言模型；糖尿病管理；LangGraph；检索增强生成；适老化设计

\newpage

\chapter*{Abstract}
\addcontentsline{toc}{chapter}{Abstract}

\textcolor{red}{[TODO: 全文完成后撰写英文摘要]}

Facing the increasingly severe aging population and chronic disease management pressure in China, community Type 2 diabetes management confronts the structural contradiction of ``limited human resources versus unlimited management demands''. This thesis designs and implements an intelligent follow-up management platform for community elderly diabetes based on Multi-Agent System (MAS). The system adopts LangGraph as the multi-agent orchestration framework and Large Language Model (LLM) as the intelligent enhancement engine, employing a hybrid architecture of ``rule-driven + LLM-enhanced''. It constructs a five-agent collaborative system including Patient Agent, Triage Agent, Scheduler Agent, Medication Agent, and Doctor Agent. Through Retrieval-Augmented Generation (RAG) technology integrating authoritative knowledge such as the ``Chinese Guidelines for Diabetes Prevention and Treatment (2024 Edition)'', the system achieves core functions including intelligent health data collection, three-color risk warning, dynamic follow-up scheduling, and medication adherence management. The front-end adopts age-friendly design with voice interaction support, lowering the usage barrier for elderly patients.

\vspace{1em}
\noindent\textbf{Keywords:} Multi-Agent System; Large Language Model; Diabetes Management; LangGraph; Retrieval-Augmented Generation; Age-friendly Design
